\documentclass[11pt,a4paper]{article}

\usepackage[margin=2.3cm]{geometry}
\usepackage{amsmath,amssymb,amsthm}
\usepackage{booktabs}
\usepackage{xcolor}
\usepackage{enumitem}
\usepackage[colorlinks=true,linkcolor=black,urlcolor=black]{hyperref}
\usepackage{parskip}
\usepackage{titlesec}
\usepackage{fancyhdr}
\usepackage{mdframed}

\definecolor{accent}{HTML}{1F4E79}
\definecolor{shade}{HTML}{F2F5F8}
\definecolor{warn}{HTML}{8B2500}
\definecolor{warnshade}{HTML}{FBF3F0}

\titleformat{\section}{\Large\bfseries\color{accent}}{\thesection.}{0.6em}{}
\titleformat{\subsection}{\large\bfseries}{}{0em}{}
\titlespacing*{\section}{0pt}{1.5em}{0.5em}

\pagestyle{fancy}
\fancyhf{}
\lhead{\small\color{accent}Buffer allocation --- the defensible core}
\rhead{\small\thepage}
\renewcommand{\headrulewidth}{0.4pt}

\newcommand{\code}[1]{\texttt{\small #1}}

\newmdenv[backgroundcolor=shade,linecolor=accent!40,linewidth=0.8pt,
          innertopmargin=8pt,innerbottommargin=8pt,skipabove=10pt,skipbelow=10pt]{keybox}
\newmdenv[backgroundcolor=warnshade,linecolor=warn!50,linewidth=0.8pt,
          innertopmargin=8pt,innerbottommargin=8pt,skipabove=10pt,skipbelow=10pt]{warnbox}

\begin{document}

\begin{center}
{\LARGE\bfseries\color{accent} Buffer allocation --- the defensible core}\\[0.4em]
{\large What to be able to say, and why each part is true}\\[0.8em]
{\small RailCast \quad$\cdot$\quad 8 September 2026}
\end{center}

\vspace{0.3em}
\begin{keybox}
\textbf{What this document is for.} The three-part argument in section 3 is the whole proof
that this is a formulation rather than a library call. It is what an interviewer will probe,
and it is the part that must be reproducible on a whiteboard without notes. Everything else
here is context for it.
\end{keybox}

\section{The claim, in one paragraph}

A timetable contains padding: scheduled running times exceed the minimum a train can
actually achieve. That padding absorbs delay. It is currently distributed however the
timetable happens to distribute it. \emph{Given the same total padding}, is there a better
allocation --- one that reduces expected lateness?

This is posed as a linear program: the padding on each segment is a decision variable,
the total is a budget constraint, delay propagation along the route is a set of linear
constraints, and expected lateness across historical days is the objective. It is solved
with HiGHS via \code{scipy.optimize.linprog}. The formulation is derived below; the solver
is only the arithmetic.

\section{The program}

Stops $i = 0,1,\dots,n$ along one train's route; segment $i$ runs from stop $i-1$ to stop
$i$. Scenarios $\omega \in \Omega$ are historical days.

\subsection*{Data}

\[
\begin{aligned}
m_i &= \text{minimum technical running time on segment } i \\
r_i &= \text{scheduled running time on segment } i \\
b_i^0 &= r_i - m_i \\
B &= \textstyle\sum_{i=1}^{n} b_i^0 \\
\delta_i^\omega &= \max\!\left(0,\ \rho_i^\omega - m_i\right)
\end{aligned}
\]

$\rho_i^\omega$ is the observed running time on segment $i$ on day $\omega$; $m_i$ is
estimated as the 5th percentile of observed running times; $b^0$ is the timetable's own
allocation and is the baseline; $B$ is the budget, so the same total is redistributed and
journey time is unchanged; $\delta_i^\omega$ is the primary delay actually incurred.

\subsection*{Variables, objective, constraints}

\[
b_i \ge 0 \quad (i = 1 \dots n), \qquad L_i^\omega \ge 0 \quad (i = 1 \dots n,\ \omega \in \Omega)
\]

\[
\min_{b,\,L}\quad \frac{1}{|\Omega|}\sum_{\omega \in \Omega}\ \sum_{i=1}^{n} w_i\, L_i^\omega
\]

\[
\begin{aligned}
\text{(C1)}\quad L_i^\omega &\ \ge\ L_{i-1}^\omega + \delta_i^\omega - b_i && \forall i,\ \forall \omega \\
\text{(C2)}\quad L_i^\omega &\ \ge\ 0 && \forall i,\ \forall \omega \\
\text{(C3)}\quad \textstyle\sum_{i=1}^{n} b_i &\ \le\ B \\
\text{(C4)}\quad 0 \le b_i &\ \le\ \bar b_i && \forall i \\
\text{(C5)}\quad L_0^\omega &\ =\ \ell_0^\omega
\end{aligned}
\]

\section{Why the maximum can be replaced by two inequalities}

This is the argument to know cold.

\subsection*{The problem}

Delay propagation along a route is genuinely nonlinear:

\[
L_i^\omega \;=\; \max\!\bigl(0,\ L_{i-1}^\omega + \delta_i^\omega - b_i\bigr).
\]

A train carries its delay forward; buffer eats into it; and it cannot become
``negatively late'' in a way that lets it depart a scheduled stop early. That $\max$ cannot
appear in a linear program.

Constraints (C1) and (C2) are the same expression written as two \emph{inequalities}. They
say $L_i^\omega$ is at least each of the two arguments --- which is weaker than saying it
equals their maximum. Every feasible point of the true system is feasible here, but so are
points where $L_i^\omega$ is too large. The relaxation must be shown to be \textbf{tight}:
that the optimum of the relaxed program is a solution of the true recursion.

\subsection*{The argument, in three steps}

\begin{keybox}
\textbf{1. The objective pushes every $L_i^\omega$ down.} Each enters with coefficient
$w_i/|\Omega| \ge 0$. Nothing in the program rewards a larger $L_i^\omega$.

\medskip
\textbf{2. Its only lower bounds are (C1) and (C2).} No other constraint mentions
$L_i^\omega$ from below. Driven down as far as feasibility allows, it comes to rest exactly
at $\max\bigl(0,\ L_{i-1}^\omega + \delta_i^\omega - b_i\bigr)$ --- the value the true
recursion specifies.

\medskip
\textbf{3. Reducing it never costs anything elsewhere.} $L_i^\omega$ appears once more, on
the right-hand side of (C1) for index $i+1$, with coefficient $+1$. Reducing $L_i^\omega$
\emph{lowers} that bound, which \emph{relaxes} the next constraint. Both forces act in the
same direction. There is no configuration in which inflating $L_i^\omega$ buys a saving.
\end{keybox}

\subsection*{The formal version}

\begin{quote}
\textbf{Claim.} Let $(b^\ast, L^\ast)$ be optimal. Then
$L_i^{\ast\omega} = \max\bigl(0,\ L_{i-1}^{\ast\omega} + \delta_i^\omega - b_i^\ast\bigr)$
for every $i,\omega$.

\textbf{Proof.} Suppose some $L_i^{\ast\omega}$ strictly exceeds that maximum. Lower it to
the maximum. Then (C1) at index $i$ still holds, since the value still meets the bound;
(C2) still holds, since the maximum is non-negative; (C1) at index $i+1$ has a strictly
\emph{smaller} right-hand side and so still holds; and (C3), (C4) are untouched, involving
only $b$. The point remains feasible, and the objective has strictly decreased ---
contradicting optimality. \hfill$\square$
\end{quote}

\subsection*{The gap in that proof, and the condition that closes it}

The last step needs the objective to \emph{strictly} decrease, which needs $w_i > 0$. Under
terminus-only weighting ($w_i = 0$ for every $i < n$) that fails, and the naive argument
does not go through.

The relaxation is still tight, for a subtler reason. Lowering $L_i^\omega$ relaxes (C1) at
$i+1$, which lets $L_{i+1}^\omega$ fall, which relaxes (C1) at $i+2$, and so on: a reduction
propagates forward along the chain with coefficient $1$ at every step. So it reaches the
terminus, where the weight is positive, and the objective does strictly decrease.

\begin{keybox}
\textbf{The general condition.} The relaxation is tight at $(i,\omega)$ provided
\[
w_j \ge 0 \ \ \text{for all } j, \qquad\text{and}\qquad \sum_{j \ge i} w_j > 0 .
\]
That is: every stop from $i$ onward has non-negative weight, and at least one of them is
strictly positive. Under uniform weighting every index qualifies immediately. Under
terminus-only weighting every index qualifies because $w_n > 0$ and the reduction
propagates forward to it.
\end{keybox}

\noindent This is the version to give if pushed. Stating only ``the objective pushes it
down'' invites the follow-up \emph{``what if that stop has zero weight?''}, and terminus-only
weighting --- which this project runs --- is exactly that case.

\section{The failure mode}

\begin{warnbox}
\textbf{Introduce a negative weight anywhere and the exactness breaks.}

Suppose $w_k < 0$ at some stop, intended to reward early arrival. Then the objective
\emph{rewards} a larger $L_k^\omega$. The solver will inflate it --- feasibly, since (C1)
and (C2) are only lower bounds --- to collect that reward, and the constraint at $k+1$
absorbs the inflation by raising $L_{k+1}^\omega$, which may be cheap if the weights beyond
$k$ are small.

The lateness variables then no longer represent lateness. The reported objective is not the
expected lateness of any timetable, and the buffer vector is optimal for a problem nobody
posed.

\medskip
\textbf{The fix is not a negative weight.} Rewarding earliness requires representing it as a
separate non-negative variable --- $E_i^\omega \ge \text{(scheduled)} - \text{(actual)}$,
carrying its own non-negative cost --- so that both quantities are minimised and neither is
inflated for gain.
\end{warnbox}

\section{Three methodological decisions}

\subsection*{1. Significance is a paired bootstrap over \emph{days}, not over rows}

The plan says to evaluate against $b^0$ on held-out days. It does not say how to decide the
difference is real, and with a handful of days an improvement of a few seconds is easily
noise. The project already has the apparatus: the paired bootstrap veto used for model
promotion.

One refinement matters. That test resamples \emph{rows}. Here the independent unit is the
\textbf{day}, not the stop-arrival: within one day, delay at consecutive stops is strongly
correlated by construction --- that is the propagation the whole model describes.
Resampling stop-arrivals would treat correlated observations as independent and understate
the variance, producing a confident answer that is not warranted.

\begin{keybox}
\textbf{The number that governs the statistics is the number of days, not the number of
variables.} A 29-stop service with 32 complete days gives 896 lateness variables and
\textbf{32} independent scenarios. Power comes from 32. Resample whole days, with
replacement, and report the interval on the difference $\big(\text{cost}(b^\ast) -
\text{cost}(b^0)\big)$.
\end{keybox}

\noindent Consequence to accept in advance: 32 scenarios is a small sample for a
significance claim, and the honest outcome may be ``an improvement was measured and cannot
be distinguished from noise.'' The way to raise power is more days, and the data source
supports it --- the API serves history back to 2007, so the scenario set can be extended by
backfilling more dates rather than by weakening the test.

\subsection*{2. The cap $\bar b_i$ needs defining, and reporting}

$\bar b_i$ appears in (C4) and is nowhere defined. Two separate questions:

\textbf{Is it needed mathematically?} No. Buffer is scarce --- (C3) binds --- so the solver
has no reason to place buffer where it cannot be used. Allocating more than
$\max_\omega\bigl(L_{i-1}^\omega + \delta_i^\omega\bigr)$ to segment $i$ is provably wasted,
because beyond that point the segment already absorbs everything reaching it in every
scenario, and the budget would be better spent elsewhere.

\textbf{Is it needed operationally?} Yes, if the answer is meant to be implementable.
Without it the optimum may schedule a four-minute segment at eleven minutes, which conflicts
with other traffic and is not a timetable anyone would adopt. A defensible choice is a
proportional cap, $\bar b_i = \alpha \, m_i$, in the spirit of the running-time supplements
timetable planners already use.

\begin{keybox}
\textbf{Report whether the cap binds.} If it does not, it is inert and the result is
data-driven. If it does, the result is partly driven by a chosen constant and that must be
stated, along with the sensitivity to $\alpha$. Run once uncapped as the reference.
\end{keybox}

\subsection*{3. Report both weightings; do not choose one}

$w_i$ is a modelling decision and boardings are unavailable. Choosing a plausible-looking
weight vector would reintroduce precisely the objection that ruled out the delay-management
formulation: an objective function containing invented numbers.

Run both defensible extremes and publish both:

\begin{center}
\begin{tabular}{@{}lll@{}}
\toprule
\textbf{Weighting} & \textbf{$w_i$} & \textbf{Optimises for} \\
\midrule
Terminus-only & $w_n = 1$, else $0$ & end-to-end punctuality \\
Uniform & $w_i = 1$ for all $i$ & passengers alighting anywhere \\
\bottomrule
\end{tabular}
\end{center}

\noindent They bracket the answer. If the optimised allocation beats $b^0$ under both, the
conclusion is robust to the weighting --- a stronger claim than either result alone. If it
wins under one and loses under the other, that is a finding about where the gain comes
from, and it is reported as such rather than resolved by picking the flattering one.

\section{Questions to expect}

\begin{enumerate}[leftmargin=1.5em,itemsep=0.3em]
  \item \emph{How is this optimization rather than calling a library?} --- Answer with the
        formulation: variables, objective, constraints, and the linearisation of the $\max$.
        The solver is arithmetic; the modelling is section 3.
  \item \emph{Your constraints are inequalities but the dynamics are an equality. Why is
        that valid?} --- The three-step argument, then the $\sum_{j\ge i} w_j > 0$ condition.
  \item \emph{What breaks it?} --- A negative weight. Explain the inflation mechanism.
  \item \emph{Why is this a linear program and not an integer program?} --- Every variable is
        continuous; there is no discrete decision. Delay management would be an integer
        program, because hold-or-drop is binary --- and it was rejected for lack of data,
        not for difficulty.
  \item \emph{How do you know the improvement is real?} --- Paired bootstrap over days, and
        the sample size is days.
  \item \emph{What are its assumptions?} --- Primary delays exogenous (the largest threat:
        drivers pace to the timetable, so added buffer may attract a slower run); minimum
        running time estimated from a sample quantile; single train, no network capacity;
        counterfactual, so simulated rather than validated.
  \item \emph{Does it assume segments are independent?} --- No. Scenarios are whole days, so
        within-day correlation is preserved exactly. This is a genuine advantage of sample
        average approximation over a parametric model.
\end{enumerate}

\vspace{0.8em}
\hrule
\vspace{0.5em}
\noindent\small\textit{If the optimiser does not beat the existing timetable, that result is
recorded and kept. A negative result honestly measured is worth more than a claim that has
to be hedged.}

\end{document}
